\documentclass[aspectratio=169, 10pt]{beamer}
\usetheme{metropolis}

% --- Edinburgh colours ----------------------------------------
\definecolor{UoERed}{HTML}{7A2318}
\definecolor{UoEGold}{HTML}{B8860B}
\definecolor{CodeBG}{HTML}{F7F3ED}
\definecolor{CodeFrame}{HTML}{D5C9B1}
\definecolor{UoEBlue}{HTML}{2a78d6}

\setbeamercolor{palette primary}{bg=UoERed, fg=white}
\setbeamercolor{frametitle}{bg=UoERed, fg=white}
\setbeamercolor{title separator}{fg=UoEGold}
\setbeamercolor{progress bar}{fg=UoEGold, bg=UoEGold!20}
\setbeamercolor{alerted text}{fg=UoERed}
\setbeamercolor{block title}{bg=UoERed!15, fg=UoERed}
\setbeamercolor{block body}{bg=UoERed!5}

% --- Packages -------------------------------------------------
\usepackage[T1]{fontenc}
\usepackage{lmodern}
\usepackage{amsmath, amssymb, mathtools}
\usepackage{listings}
\usepackage{graphicx, booktabs}
\usepackage{tikz}
\usetikzlibrary{arrows.meta, positioning, calc, decorations.pathreplacing}
\usepackage{hyperref}
\hypersetup{colorlinks=true, linkcolor=UoERed, urlcolor=UoERed!70!black}
\setbeamertemplate{navigation symbols}{}

\lstset{
  language=Python,
  basicstyle=\ttfamily\scriptsize,
  keywordstyle=\color{UoERed}\bfseries,
  commentstyle=\color{gray}\itshape,
  stringstyle=\color{UoEGold},
  backgroundcolor=\color{CodeBG},
  frame=single,
  rulecolor=\color{CodeFrame},
  numbers=none, breaklines=true, showstringspaces=false, tabsize=4,
  xleftmargin=4pt, xrightmargin=4pt, aboveskip=6pt, belowskip=4pt,
}

\AtBeginSection[]{
  \begin{frame}
    \vfill\centering
    \usebeamerfont{title}\insertsectionhead\par
    \vfill
  \end{frame}
}

\title{Week 1 --- Introduction to Forecasting}
\subtitle{Time Series Components, Stationarity, and Naive Methods}
\author{Dr Juan Zurita}
\institute{Economic \& Business Forecasting\\
The University of Edinburgh~$\cdot$~School of Economics}
\date{}

\begin{document}
\maketitle


\section{What Is Forecasting?}

\begin{frame}{Why Forecast?}
\begin{itemize}
    \item \alert{Monetary policy}: central banks forecast inflation to set interest rates
    \item \alert{Fiscal planning}: governments forecast GDP to plan budgets
    \item \alert{Business}: firms forecast demand for inventory and staffing
    \item \alert{Finance}: risk models rely on forecasts of returns and volatility
\end{itemize}

\vspace{0.5em}
\textbf{Key question:} given observations $y_1, y_2, \ldots, y_T$, what is our best prediction of $y_{T+h}$ for horizon $h \geq 1$?
\end{frame}

\begin{frame}{Types of Forecasting}
\begin{columns}[T]
\begin{column}{0.48\textwidth}
\textbf{Qualitative}
\begin{itemize}
    \item Expert judgement
    \item Delphi method
    \item Surveys
\end{itemize}
\end{column}
\begin{column}{0.48\textwidth}
\textbf{Quantitative}
\begin{itemize}
    \item Time-series models
    \item Causal / econometric models
    \item Machine learning
\end{itemize}
\end{column}
\end{columns}
\vspace{1em}
This course focuses on \alert{quantitative methods}.
\end{frame}

\section{Time Series Components}

\begin{frame}{Decomposition}
Any time series can be decomposed:
\[
y_t = T_t + S_t + C_t + \varepsilon_t
\]
\begin{description}
    \item[$T_t$] \textbf{Trend} --- long-run increase or decrease
    \item[$S_t$] \textbf{Seasonality} --- regular periodic pattern (fixed period)
    \item[$C_t$] \textbf{Cycle} --- fluctuations without a fixed period
    \item[$\varepsilon_t$] \textbf{Irregular} --- random noise
\end{description}
\end{frame}

\begin{frame}{Additive vs Multiplicative}
\begin{columns}[T]
\begin{column}{0.48\textwidth}
\textbf{Additive:}
\[
y_t = T_t + S_t + \varepsilon_t
\]
Seasonal amplitude is \emph{constant}.
\end{column}
\begin{column}{0.48\textwidth}
\textbf{Multiplicative:}
\[
y_t = T_t \times S_t \times \varepsilon_t
\]
Seasonal amplitude \emph{grows} with the level.
\end{column}
\end{columns}
\end{frame}

\section{Stationarity}

\begin{frame}{Covariance Stationarity}
A time series $\{y_t\}$ is \alert{covariance stationary} if:
\begin{enumerate}
    \item $E[y_t] = \mu$ for all $t$ (constant mean)
    \item $\operatorname{Var}(y_t) = \sigma^2$ for all $t$ (constant variance)
    \item $\operatorname{Cov}(y_t, y_{t-k}) = \gamma_k$ depends only on $k$, not $t$
\end{enumerate}

\vspace{0.5em}
\textbf{Why it matters:} most forecasting models assume stationarity, or that we can transform the data to achieve it.
\end{frame}

\begin{frame}{Testing for Stationarity}
\begin{itemize}
    \item \textbf{Visual inspection}: does the series wander or mean-revert?
    \item \textbf{ACF}: slow decay $\Rightarrow$ likely non-stationary
    \item \textbf{Augmented Dickey-Fuller (ADF) test}: $H_0$: unit root (non-stationary)
    \item \textbf{KPSS test}: $H_0$: stationary
\end{itemize}

\vspace{0.5em}
\textbf{Common remedy:} differencing --- $\Delta y_t = y_t - y_{t-1}$.
\end{frame}

\section{Autocorrelation}

\begin{frame}{ACF and PACF}
\begin{columns}[T]
\begin{column}{0.48\textwidth}
\textbf{Autocorrelation Function (ACF):}
\[
\rho_k = \frac{\gamma_k}{\gamma_0} = \frac{\operatorname{Cov}(y_t, y_{t-k})}{\operatorname{Var}(y_t)}
\]
Measures \emph{total} correlation at lag $k$.
\end{column}
\begin{column}{0.48\textwidth}
\textbf{Partial ACF (PACF):}

Correlation between $y_t$ and $y_{t-k}$ \emph{after} removing the effect of lags $1, \ldots, k-1$.
\end{column}
\end{columns}

\vspace{0.5em}
ACF and PACF are the main tools for \alert{model identification}.
\end{frame}

\section{A First Forecast}

\begin{frame}{The Naive Method}
\textbf{Forecast:} $\hat{y}_{T+h|T} = y_T$ \quad (repeat the last observation)

\vspace{0.5em}
\begin{itemize}
    \item Surprisingly hard to beat for \alert{random walks}
    \item Serves as a \textbf{benchmark} --- any useful model should beat it
    \item Seasonal variant: $\hat{y}_{T+h|T} = y_{T+h-m}$ (same season last year)
\end{itemize}

\vspace{0.5em}
\textbf{Forecast evaluation:}
\[
\text{RMSE} = \sqrt{\frac{1}{H}\sum_{h=1}^{H}\left(y_{T+h} - \hat{y}_{T+h|T}\right)^2}
\]
\end{frame}

\begin{frame}{Summary}
\begin{itemize}
    \item Time series have \textbf{trend}, \textbf{seasonal}, \textbf{cyclical}, and \textbf{irregular} components
    \item \textbf{Stationarity} is the key assumption behind most models
    \item \textbf{ACF/PACF} reveal the dependence structure
    \item The \textbf{naive method} is a useful benchmark
\end{itemize}

\vspace{0.5em}
\textbf{Next week:} probability and statistics foundations for forecasting.
\end{frame}


\end{document}
